\documentclass[UTF8,zihao=-4,oneside,fontset=fandol]{ctexrep}

\usepackage[a4paper,margin=2.35cm]{geometry}
\usepackage{amsmath,amssymb,amsthm,mathtools,bm}
\usepackage{booktabs,tabularx,array,multirow}
\usepackage{enumitem}
\usepackage{xcolor}
\usepackage{hyperref}
\usepackage{fancyhdr}
\usepackage{setspace}
\usepackage{listings}
\usepackage{tikz}
\usepackage[most]{tcolorbox}
\usetikzlibrary{arrows.meta,positioning,shapes.geometric,calc}

\definecolor{DeepBlue}{HTML}{1F4E79}
\definecolor{DeepGreen}{HTML}{2E6B4F}
\definecolor{WarmRed}{HTML}{A33A32}
\definecolor{SoftGray}{HTML}{F4F5F6}

\hypersetup{
  colorlinks=true,
  linkcolor=DeepBlue,
  urlcolor=DeepBlue,
  citecolor=DeepGreen,
  pdftitle={Lean 4 与 mathlib 自学讲义},
  pdfauthor={RUCmonk2}
}

\setstretch{1.22}
\setlength{\parindent}{2em}
\setlength{\parskip}{0.25em}
\setlist{itemsep=0.25em,topsep=0.35em}

\pagestyle{fancy}
\fancyhf{}
\fancyhead[L]{Lean 4 与 mathlib 自学讲义}
\fancyhead[R]{\leftmark}
\fancyfoot[C]{\thepage}
\setlength{\headheight}{14pt}

\newtheorem{definition}{定义}[chapter]
\newtheorem{theorem}[definition]{定理}
\newtheorem{proposition}[definition]{命题}
\newtheorem{lemma}[definition]{引理}
\newtheorem{corollary}[definition]{推论}
\theoremstyle{remark}
\newtheorem{example}[definition]{例}
\newtheorem{remark}[definition]{说明}

\newtcolorbox{intuitionbox}{
  colback=DeepBlue!4,colframe=DeepBlue!70!black,
  title=直觉,breakable,boxrule=0.7pt,arc=1mm
}
\newtcolorbox{warningbox}{
  colback=WarmRed!4,colframe=WarmRed!75!black,
  title=容易混淆,breakable,boxrule=0.7pt,arc=1mm
}
\newtcolorbox{aibox}{
  colback=DeepGreen!4,colframe=DeepGreen!75!black,
  title=与 AI4Math / AI4TCS 的连接,breakable,boxrule=0.7pt,arc=1mm
}
\newtcolorbox{checkpointbox}{
  colback=SoftGray,colframe=black!50,
  title=本章验收,breakable,boxrule=0.6pt,arc=1mm
}

\lstdefinelanguage{Lean}{
  morekeywords={abbrev,axiom,by,class,def,deriving,do,else,end,example,
    extends,for,forall,fun,if,import,inductive,instance,let,match,
    namespace,open,opaque,partial,private,protected,return,structure,
    theorem,then,universe,variable,where,with},
  sensitive=true,
  morecomment=[l]{--},
  morecomment=[s]{/-}{-/},
  morestring=[b]"
}
\lstset{
  basicstyle=\ttfamily\small,
  breaklines=true,
  columns=fullflexible,
  frame=single,
  rulecolor=\color{black!25},
  backgroundcolor=\color{black!2},
  xleftmargin=0.6em,
  xrightmargin=0.6em,
  showstringspaces=false,
  language=Lean
}

\newcommand{\code}[1]{\texttt{#1}}
\newcommand{\Nat}{\mathbb{N}}
\newcommand{\PropT}{\mathsf{Prop}}
\newcommand{\TypeT}{\mathsf{Type}}

\title{\heiti Lean 4 与 mathlib 自学讲义\\[0.5em]
\Large 从可执行定义、小型证明到可审计 verifier}
\author{RUCmonk2}
\date{2026 年 8 月}

\begin{document}

\hypersetup{pageanchor=false}
\maketitle
\hypersetup{pageanchor=true}

\chapter*{如何使用这本讲义}
\addcontentsline{toc}{chapter}{如何使用这本讲义}

本讲义面向第一次系统学习 Lean 4、同时希望把它用于 AI4Math、AI4TCS 和随机算法形式化的读者。目标不是背 tactic，而是建立一条稳定链路：把自然语言命题写成精确类型，读懂 proof state，构造能被 kernel 检查的证明项，并诚实记录尚未形式化的语义与性能性质。

建议每个例子都完成四次检查：
\begin{enumerate}
  \item 在纸上写出对象、假设和结论，不从 tactic 开始；
  \item 用 \code{\#check} 确认关键表达式的类型；
  \item 在固定 toolchain 的 Lake 项目中编译，而不只依赖编辑器提示；
  \item 故意删掉一个假设或改错一个类型，观察 Lean 的错误信息如何变化。
\end{enumerate}

代码示例以 Lean 4 与 mathlib 的常见写法为主。mathlib API 会随版本演进，学习项目必须提交 \code{lean-toolchain} 与 Lake 配置；当示例和本机版本不一致时，应先检索当前版本的定义和定理，而不是盲目改到“没有红线”。

\tableofcontents
\clearpage

\chapter{环境、Lake 项目与内核检查模型}

\section{Lean 工具链由什么组成}

第一次接触 Lean 时，容易把 VS Code 中的彩色目标窗口当成 Lean 本身。实际上常用工作流至少包含四层：
\begin{description}
  \item[elan] 安装并切换 Lean toolchain，作用类似 Rust 的 \code{rustup}；
  \item[Lean] 解析、展开和检查 Lean 源文件，最终由小型 kernel 验证证明项；
  \item[Lake] 管理项目、依赖和构建任务；
  \item[编辑器扩展] 把当前位置的目标、诊断和补全展示出来，但它不是最终可信边界。
\end{description}

先在终端记录真实版本：
\begin{lstlisting}[language={}]
elan --version
lean --version
lake --version
\end{lstlisting}
同一个项目应固定 \code{lean-toolchain}，复制代码时也要带上 toolchain 与依赖版本。缺少这些信息时，“在他那里能编译”不能推出“在当前环境能编译”。

\section{创建最小 Lake 项目}

在空目录中可以运行：
\begin{lstlisting}[language={}]
lake init LeanSelfStudy
lake build
\end{lstlisting}
不同 Lean 版本生成的模板可能略有差别，但应至少识别以下对象：
\begin{center}
\begin{tabularx}{0.94\textwidth}{>{\bfseries}lX}
\toprule
文件或目录 & 作用 \\
\midrule
\code{lean-toolchain} & 固定 Lean 版本；构建复现的第一入口 \\
\code{lakefile.toml} 或 \code{lakefile.lean} & 声明包、库、可执行程序和依赖 \\
\code{Main.lean} & 模板中的可执行入口，若项目需要命令行程序 \\
\code{LeanSelfStudy/} & 自己的库模块；文件路径通常对应模块名 \\
\code{lake-manifest.json} & Lake 解析后的依赖锁定信息 \\
\code{.lake/} & 本机下载与构建缓存，不应当作学习成果 \\
\bottomrule
\end{tabularx}
\end{center}

引入 mathlib 时，应按 mathlib 当前文档创建项目或添加依赖，再运行 \code{lake update} 和 \code{lake build}。不要手工复制 mathlib 的缓存目录；项目配置和 manifest 才是可审计依赖。

\section{第一个可检查文件}

把下面内容放进项目模块中：
\begin{lstlisting}
import Mathlib

#check Nat
#check Nat.add
#check Nat.add_comm

def double (n : Nat) : Nat := n + n

#eval double 6

example (n : Nat) : double n = n + n := by
  rfl
\end{lstlisting}

\code{\#check} 让 elaborator 告诉我们表达式的类型；\code{\#eval} 对可执行定义求值；\code{example} 声明一个无需命名保存的定理。最后的 \code{rfl} 能成功，是因为展开 \code{double} 后两边定义性相同。

单独检查文件和检查整个项目分别使用：
\begin{lstlisting}[language={}]
lake env lean LeanSelfStudy/Basic.lean
lake build
\end{lstlisting}
前者适合缩小错误，后者负责确认模块依赖和整个项目都一致。

\section{从源码到 kernel 的路径}

Lean 不直接把你键入的 tactic 字符串交给 kernel。可以用下图理解处理链：
\begin{center}
\begin{tikzpicture}[
  node distance=7mm,
  box/.style={draw,rounded corners=1mm,align=center,minimum width=3.2cm,minimum height=8mm},
  arr/.style={-{Stealth[length=2mm]},thick}
]
\node[box] (src) {源码与 tactic};
\node[box,below=of src] (elab) {parser / elaborator\\补全隐式参数与实例};
\node[box,below=of elab] (term) {核心表达式与证明项};
\node[box,below=of term] (kernel) {kernel 类型检查};
\node[box,below=of kernel] (ok) {接受或拒绝};
\draw[arr] (src)--(elab);
\draw[arr] (elab)--(term);
\draw[arr] (term)--(kernel);
\draw[arr] (kernel)--(ok);
\end{tikzpicture}
\end{center}

Tactic、搜索器甚至语言模型都可以帮助构造证明项；最终 kernel 只检查这个项是否具有声明的类型。由此得到两个同时成立的判断：
\begin{enumerate}
  \item 自动化很强不必然扩大可信计算基；只要最终产生普通证明项，kernel 仍会复查；
  \item kernel 接受只保证“形式语句被证明”，不保证自然语言题目被正确翻译，也不保证定理有研究意义。
\end{enumerate}

\begin{aibox}
AI proof agent 可以生成 statement、tactic 或完整 proof term，Lean 负责形式正确性。人类仍需检查 statement fidelity、依赖中是否引入过强公理、问题是否被偷换，以及算法复杂度等未进入命题的性质。
\end{aibox}

\section{错误信息先按层分类}

面对失败时，先问它发生在哪一层：
\begin{description}
  \item[解析错误] 括号、缩进、关键字或 Unicode/ASCII 记号不合法；
  \item[名称解析错误] 没有导入模块、命名空间不对或 API 已改名；
  \item[elaboration 错误] 隐式参数、类型类实例或预期类型无法确定；
  \item[证明目标未关闭] tactic 执行过，但仍有 goals；
  \item[构建错误] 单文件似乎可用，但模块名、依赖或项目配置不一致。
\end{description}

\begin{warningbox}
重启编辑器、删除全部缓存偶尔能解决环境问题，却不能替代错误分类。若一个 10 行例子仍然失败，应先保存完整错误和 proof state，再最小化代码；否则一次“偶然变绿”不会留下可复用理解。
\end{warningbox}

\section{最小可复现环境记录}

一个可以交给别人复查的 Lean 结果至少应包含：
\begin{enumerate}
  \item Lean 与 mathlib 的固定版本；
  \item Lake 配置、manifest 和所有必要源文件；
  \item 从干净 checkout 开始的构建命令；
  \item 定理是否使用 \code{axiom}、\code{sorry} 或非标准代码生成路径；
  \item 失败时的最小例子和完整错误，而不是只有截图。
\end{enumerate}

\begin{checkpointbox}
\begin{enumerate}
  \item 解释 elan、Lean、Lake、编辑器扩展各自负责什么；
  \item 创建一个固定 toolchain 的项目，并用终端完成一次 \code{lake build}；
  \item 说明 tactic、elaborator、proof term 与 kernel 的关系；
  \item 把一次失败归类为解析、名称、elaboration、证明或构建错误；
  \item 说清 Lean kernel 接受一条定理能保证什么、不能保证什么。
\end{enumerate}
\end{checkpointbox}

\chapter{表达式、定义、函数与计算}

\section{每个表达式都有类型}

Lean 的核心判断写成
\[
\Gamma \vdash t : T,
\]
表示在局部环境 $\Gamma$ 中，项 $t$ 具有类型 $T$。例如自然数 \code{3} 的类型是 \code{Nat}，函数 \code{Nat.succ} 的类型是 \code{Nat -> Nat}，而一个等式证明也是某个命题类型的项。

\begin{lstlisting}
#check 3
#check Nat.succ
#check fun n : Nat => n + 1
#check Type
#check Prop
\end{lstlisting}

类型本身也需要类型。Lean 使用 universe 层级避免“所有类型组成的类型包含自身”带来的悖论。入门时可以先记住：\code{Nat : Type}，而 \code{Type : Type 1}；多态定义经常让 universe 参数由 elaborator 自动补全。

\section{函数类型与函数定义}

若输入类型为 $A$、输出类型为 $B$，函数类型写作 $A\to B$。匿名函数和命名定义分别写作：
\begin{lstlisting}
#check fun n : Nat => n * n

def square (n : Nat) : Nat := n * n

def compose (f : B -> C) (g : A -> B) : A -> C :=
  fun x => f (g x)
\end{lstlisting}
最后一个定义中的 \code{A B C} 需要预先作为类型变量声明：
\begin{lstlisting}
universe u v w
variable {A : Type u} {B : Type v} {C : Type w}
\end{lstlisting}

函数默认柯里化。\code{Nat -> Nat -> Nat} 解析为 \code{Nat -> (Nat -> Nat)}：先接受一个自然数，再返回一个接受自然数的函数。因此多参数应用写作 \code{f x y}，而不是必须构造元组。

\section{显式参数、隐式参数与命名参数}

圆括号参数必须显式提供，花括号参数通常由上下文推断：
\begin{lstlisting}
def identity {A : Type} (x : A) : A := x

#check identity 5
#check identity "lean"
#check @identity
\end{lstlisting}
\code{@identity} 要求显示包括隐式参数在内的完整函数。诊断类型推断时，还可以用命名参数指定：\code{identity (A := Nat) 5}。

\begin{warningbox}
“参数没有写出来”不等于“参数不存在”。mathlib 定理常含 universe、类型、实例和假设等多层隐式参数。应用失败时先 \code{\#check} 定理，再必要时查看 \code{@theoremName} 或打开更详细的 pretty printer。
\end{warningbox}

\section{局部定义、命名空间与模块}

\code{let} 创建局部名称；\code{namespace} 避免全局命名冲突：
\begin{lstlisting}
namespace SelfStudy

def twice (f : Nat -> Nat) (x : Nat) : Nat :=
  let y := f x
  f y

#eval twice (fun n => n + 1) 10

end SelfStudy
\end{lstlisting}
跨文件使用定义时，需要让文件路径、模块名与 \code{import} 对应。\code{open} 只改变名称访问方式，并不会导入尚未加载的模块。

\section{模式匹配与递归}

代数数据类型通过构造子产生值，模式匹配负责拆解值：
\begin{lstlisting}
def isZero : Nat -> Bool
  | 0 => true
  | _ + 1 => false

def sum : List Nat -> Nat
  | [] => 0
  | x :: xs => x + sum xs

#eval sum [2, 3, 5]
\end{lstlisting}
\code{sum} 的递归调用作用于严格更小的尾列表，所以 Lean 的终止检查器可以接受。逻辑一致性要求一般递归不能无条件地产生任意类型的值；若递归不是结构递减，需要提供终止度量与证明。

\begin{intuitionbox}
函数式定义同时承担三种角色：它是可执行程序，是数学对象，也是后续证明展开的重写规则。定义接口若写得清楚，计算和证明会共享同一份结构；定义若混入过多实现细节，后续 theorem 也会被这些细节绑住。
\end{intuitionbox}

\section{归约与定义性相等}

Lean 中两类“相等”需要分开：
\begin{description}
  \item[定义性相等] 通过展开定义、函数应用归约或模式匹配计算后相同，kernel 可直接识别；
  \item[命题相等] 类型为 \code{Eq a b} 的命题，需要一个证明项。
\end{description}

例如：
\begin{lstlisting}
def addOne (n : Nat) := n + 1

example : addOne 2 = 3 := by
  rfl

example (n : Nat) : n + 0 = n := by
  simpa using Nat.add_zero n
\end{lstlisting}
第一条可通过计算完成；第二条对变量 $n$ 并不是简单把两边算成数字，需要自然数加法的定理（该定理本身由归纳证明）。

\section{可执行不等于已证明}

\code{\#eval} 能展示某个输入的计算结果，但它没有证明对所有输入成立。例如运行排序函数得到有序列表，只验证了一个样例。一般性正确性需要分别说明：输出有序、输出是输入的排列，以及函数终止。

\begin{aibox}
AI4TCS 中模型可以生成候选程序，\code{\#eval}、单元测试和随机测试用于快速筛选；Lean theorem 用于证明明确写入类型的性质。测试与证明共享定义，却提供不同等级的证据，不能互相冒充。
\end{aibox}

\section{定义设计检查表}

在开始证明前，先检查：
\begin{enumerate}
  \item 输入、输出和错误情况是否在类型中表达清楚；
  \item 使用 \code{Nat}、\code{Int}、\code{Fin n} 还是一般类型是否符合问题；
  \item 无效状态能否被构造，若能，后续需要什么不变量；
  \item 定义是面向规范还是面向高性能实现，两者是否需要分层；
  \item 递归为何终止，模式是否覆盖所有构造子。
\end{enumerate}

\begin{checkpointbox}
\begin{enumerate}
  \item 用 \code{\#check} 解释五个表达式的完整类型；
  \item 写出一个带隐式类型参数的多态函数，并用 \code{@} 查看完整参数；
  \item 定义列表乘积或最大值的安全版本，说明空列表怎样处理；
  \item 区分定义性相等与需要证明的命题相等；
  \item 解释为什么 \code{\#eval} 的十个样例不能代替全称定理。
\end{enumerate}
\end{checkpointbox}

\chapter{命题即类型与基础证明}

\section{Curry--Howard 对应}

Lean 中命题位于 \code{Prop}，证明命题 $P$ 就是构造一个类型为 $P$ 的项。常见逻辑连接词对应类型构造：
\begin{center}
\begin{tabularx}{0.94\textwidth}{>{\bfseries}l l X}
\toprule
逻辑 & Lean 形式 & 证明数据 \\
\midrule
$P\Rightarrow Q$ & \code{P -> Q} & 把 $P$ 的证明映到 $Q$ 的证明的函数 \\
$P\land Q$ & \code{And P Q} & 同时含 $P$ 与 $Q$ 证明的结构 \\
$P\lor Q$ & \code{Or P Q} & 左分支或右分支，并携带相应证明 \\
$\forall x, P(x)$ & \code{forall x, P x} & 对每个 $x$ 返回证明的依赖函数 \\
$\exists x, P(x)$ & \code{Exists fun x => P x} & 一个 witness 和它满足性质的证明 \\
$\neg P$ & \code{Not P} & 从 $P$ 的证明导出 \code{False} 的函数 \\
\bottomrule
\end{tabularx}
\end{center}

这一观点把“证明步骤”变成精确的构造操作，而不是自然语言中模糊的“显然”。

\section{直接证明与 tactic proof}

同一个定理可以直接写 proof term，也可以用 tactic 逐步构造：
\begin{lstlisting}
example (p q : Prop) (hp : p) (hq : q) : And p q :=
  And.intro hp hq

example (p q : Prop) (hp : p) (hq : q) : And p q := by
  constructor
  exact hp
  exact hq
\end{lstlisting}
第二段执行 \code{constructor} 后产生两个目标。编辑器展示的 proof state 可以抽象成：
\[
\underbrace{p,q:\PropT,\ hp:p,\ hq:q}_{\text{local context}}
\quad\vdash\quad
\underbrace{p\land q}_{\text{target}}.
\]
每个 tactic 都应当被理解为如何改变这个 judgment。

\section{蕴含、全称量词与引入}

证明函数类型时，引入输入：
\begin{lstlisting}
example (p q : Prop) : p -> q -> p := by
  intro hp
  intro hq
  exact hp

example : forall n : Nat, n = n := by
  intro n
  rfl
\end{lstlisting}
\code{intro} 不是“假设结论成立”，而是按照函数类型构造一个接受证明或对象的函数。

\section{合取、析取与分类讨论}

合取的使用是拆字段，析取的使用必须覆盖两个分支：
\begin{lstlisting}
example (p q : Prop) (h : And p q) : And q p := by
  cases h with
  | intro hp hq => exact And.intro hq hp

example (p q : Prop) (h : Or p q) : Or q p := by
  cases h with
  | inl hp => exact Or.inr hp
  | inr hq => exact Or.inl hq
\end{lstlisting}
\code{cases} 按 \code{And.intro} 构造子拆出两个字段。Lean 也支持更紧凑的结构模式，但它不会减少需要提供的数据。

\section{存在量词必须给出 witness}

\begin{lstlisting}
example : Exists fun n : Nat => n + 1 = 4 := by
  refine Exists.intro 3 ?_
  rfl
\end{lstlisting}
\code{refine} 提供带“洞”的部分证明，\code{?\_} 变成新目标。存在性证明的核心是 witness $3$，而不是 tactic 的名字。

\section{否定、矛盾与底类型}

\code{False} 没有构造子；一旦得到它，可以推出任意命题：
\begin{lstlisting}
example (p : Prop) (hp : p) (hnp : Not p) : False := by
  exact hnp hp

example (p q : Prop) (h : p -> q) (hnq : Not q) : Not p := by
  intro hp
  exact hnq (h hp)
\end{lstlisting}
第二条定理的目标 \code{Not p} 展开就是 \code{p -> False}，所以先 \code{intro hp}。

\section{等式、重写与简化}

若有 $h:a=b$，\code{rw [h]} 按方向把 $a$ 替换为 $b$；\code{rw [<- h]} 反向替换。\code{simp} 使用标记过的简化规则反复重写：
\begin{lstlisting}
example (a b c : Nat) (h : a = b) : a + c = b + c := by
  rw [h]

example (xs : List Nat) : [] ++ xs = xs := by
  simp
\end{lstlisting}
\code{rw} 更适合表达“使用了哪条等式”；\code{simp} 更适合清理单位元、构造子和已注册规则。遇到失败时，先明确期望哪一处发生哪种替换。

\section{常用 tactic 的逻辑角色}

\begin{center}
\begin{tabularx}{0.94\textwidth}{>{\ttfamily}lX}
\toprule
\normalfont tactic & 主要作用 \\
\midrule
intro & 为函数、蕴含或全称目标引入变量/假设 \\
exact & 提供类型与当前目标完全匹配的证明项 \\
apply & 反向使用定理，把当前目标变成它的前提 \\
constructor & 使用目标类型的构造子并产生字段目标 \\
cases & 按归纳类型或析取的构造子分支讨论 \\
rcases & 用模式拆解存在、合取或结构 \\
rw & 按给定等式定向重写 \\
simp & 用简化规则集做规范化重写 \\
have & 声明局部中间结论，控制证明结构 \\
show & 重新表述当前目标，帮助 elaboration 或阅读 \\
\bottomrule
\end{tabularx}
\end{center}

\section{构造逻辑与经典逻辑}

Lean 的核心逻辑是构造性的。排中律、双重否定消去等经典原理可以显式导入或局部使用，但应让依赖可见：
\begin{lstlisting}
example (p : Prop) : Or p (Not p) := by
  classical
  exact Classical.em p
\end{lstlisting}
形式化数学中使用经典逻辑通常完全合理；重要的是不要把“可构造 witness”与“经典地证明存在”混为一谈，尤其当后续希望从证明提取算法时。

\begin{warningbox}
Tactic 成功不等于你已经理解定理。至少应能回答：当前目标是什么、tactic 使用了哪个构造子或定理、产生了哪些子目标、是否引入经典公理，以及删除哪个假设后证明会失败。
\end{warningbox}

\begin{checkpointbox}
\begin{enumerate}
  \item 分别用 proof term 和 tactic 证明合取交换；
  \item 证明 \code{(p -> q) -> (Not q -> Not p)}，逐行解释 proof state；
  \item 构造一个自然数 witness，证明某个简单存在命题；
  \item 用 \code{rw} 与 \code{simp} 各完成一个等式证明，并比较两者；
  \item 找出一个需要经典逻辑的命题，明确标出 \code{classical} 的作用域。
\end{enumerate}
\end{checkpointbox}

\chapter{归纳类型、递归与归纳证明}

\section{归纳类型由构造子生成}

自然数、列表、树和语法树都可以看成“由有限组构造子生成的最小类型”。例如一个二叉树：
\begin{lstlisting}
inductive BinTree (A : Type) where
  | empty : BinTree A
  | node : BinTree A -> A -> BinTree A -> BinTree A
deriving Repr
\end{lstlisting}
这个声明同时提供：构造值的方法、按构造子分类讨论的原则，以及递归/归纳时可使用的 eliminator。任何 \code{BinTree A} 都由 \code{empty} 或 \code{node} 有限构造而来。

\section{模式匹配定义函数}

\begin{lstlisting}
namespace BinTree

def size : BinTree A -> Nat
  | .empty => 0
  | .node left _ right => size left + 1 + size right

def mirror : BinTree A -> BinTree A
  | .empty => .empty
  | .node left x right => .node (mirror right) x (mirror left)

end BinTree
\end{lstlisting}
每个分支对应一个构造子。Lean 检查模式是否覆盖所有情况，并检查递归调用是否作用于结构上更小的子项。

\begin{intuitionbox}
归纳类型给出“数据能怎样产生”，递归定义给出“怎样消费数据得到一个值”，归纳证明给出“怎样对所有可能的产生方式证明性质”。三者共享同一个构造子结构。
\end{intuitionbox}

\section{自然数归纳}

要证明 $P(n)$ 对所有 $n:\Nat$ 成立，需要证明 $P(0)$，以及从 $P(n)$ 推出 $P(n+1)$。例如：
\begin{lstlisting}
theorem add_zero_self (n : Nat) : n + 0 = n := by
  induction n with
  | zero => rfl
  | succ n ih =>
      simp [Nat.add, ih]
\end{lstlisting}
这里 \code{ih} 的类型是前一自然数上的归纳假设。实际项目通常直接使用 mathlib 已有的 \code{Nat.add\_zero}；手工证明的目的在于看清归纳原则，而不是重复造库。

\section{列表归纳}

列表只有空表与 \code{cons} 两种构造：
\begin{lstlisting}
theorem append_nil_self (xs : List A) : xs ++ [] = xs := by
  induction xs with
  | nil => rfl
  | cons x xs ih =>
      simp [ih]
\end{lstlisting}
在 \code{cons} 分支，归纳假设只覆盖尾列表 \code{xs}。若结论涉及额外参数，参数是否在归纳前引入，会影响归纳假设的强弱。

\section{树归纳与局部假设}

二叉树节点有两个递归子树，所以节点分支获得两个归纳假设：
\begin{lstlisting}
namespace BinTree

theorem size_mirror (t : BinTree A) :
    size (mirror t) = size t := by
  induction t with
  | empty => rfl
  | node left x right ihLeft ihRight =>
      simp [mirror, size, ihLeft, ihRight,
        Nat.add_assoc, Nat.add_comm, Nat.add_left_comm]

end BinTree
\end{lstlisting}
证明的数学内容是：镜像交换左右子树，但加法可交换，因此节点总数不变。\code{simp} 负责展开定义并使用归纳假设；交换律和结合律负责规范化加法顺序。

\section{先找不变量，再做归纳}

算法证明中的归纳对象不一定是一个数字，也可能是：
\begin{itemize}
  \item 输入列表或语法树；
  \item 已处理的流前缀；
  \item 状态转换的步数；
  \item 执行关系的推导树；
  \item 递归调用的结构或终止度量。
\end{itemize}

假设顺序处理输入 $x_1,\ldots,x_n$，状态满足
\[
\operatorname{Inv}(s_i,x_1,\ldots,x_i).
\]
典型证明由初始化与保持性组成：
\[
\operatorname{Inv}(s_0,[]) ,\qquad
\operatorname{Inv}(s_i,p_i)\Rightarrow
\operatorname{Inv}(\operatorname{step}(s_i,x_{i+1}),p_{i+1}).
\]
然后对前缀或步数归纳，得到所有执行阶段的不变量。

\section{归纳假设太弱怎么办}

常见症状是：进入归纳分支后，\code{ih} 只对某个固定参数成立，而目标需要它对变化后的参数成立。处理顺序应当是：
\begin{enumerate}
  \item 在纸上写出真正需要的加强命题；
  \item 决定哪些参数应在归纳前保持一般；
  \item 必要时用 \code{generalizing} 或先 \code{revert} 参数；
  \item 只有命题强度正确后，才继续尝试自动化。
\end{enumerate}

例如证明带累加器的尾递归函数正确时，通常应证明“对任意初始累加器”的更强命题，而不是只证明累加器为零的特例。

\begin{warningbox}
反复尝试 \code{simp}、\code{aesop} 不能修复错误的归纳命题。若归纳假设在数学上不足，搜索更多 tactic 只会掩盖真正问题：需要加强 theorem statement 或改变归纳对象。
\end{warningbox}

\section{良基递归与终止证明}

结构递归之外，欧几里得算法、快速排序或图搜索常按数值度量下降。Lean 允许为递归提供终止依据，例如列表长度、区间大小或一个良基关系。证明需要明确每次调用的度量严格减小。

对学习项目，建议先把函数拆成：
\begin{enumerate}
  \item 纯粹的一步转换；
  \item 显式 fuel 的有限执行器；
  \item 关于 fuel 或度量的单调/终止 lemma；
  \item 最终的部分正确性与终止性结论。
\end{enumerate}
这比一开始把复杂递归、状态和正确性揉进一个定义更容易调试。

\section{归纳证明的审计问题}

阅读任何归纳证明时都问：归纳对象是什么？构造子是否覆盖完全？每个归纳假设对应哪一个递归子对象？命题是否为了归纳而加强？证明只得到部分正确性，还是也包含终止性？

\begin{aibox}
语言模型生成归纳证明时，最常见的失败并不是不会写 tactic，而是选错归纳变量或忘记泛化参数。一个 verifier 应记录失败 proof state 和归纳假设形状，让搜索器能够修改 statement/induction schema，而不只是继续采样 tactic。
\end{aibox}

\begin{checkpointbox}
\begin{enumerate}
  \item 自定义一个带叶子值的树，并写 \code{size} 与 \code{map}；
  \item 对列表证明一个结构归纳定理，指出 base 与 step；
  \item 证明树镜像两次得到原树，标出两个归纳假设；
  \item 举出一个需要加强归纳命题的累加器例子；
  \item 区分结构递归、良基递归、部分正确性和终止性。
\end{enumerate}
\end{checkpointbox}

\chapter{依赖类型、结构、子类型与等式}

\section{返回类型可以依赖输入}

普通函数 $A\to B$ 的返回类型固定为 $B$；依赖函数的类型写作
\[
(x:A)\to B(x),
\]
返回类型会随输入 $x$ 改变。Lean 的全称量词本身就是依赖函数类型。

有限索引是最实用的例子之一：\code{Fin n} 表示小于 $n$ 的自然数。列表安全索引可以要求一个随列表长度变化的证明：
\begin{lstlisting}
def safeGet (xs : List A) (i : Fin xs.length) : A :=
  xs.get i
\end{lstlisting}
调用者不能提供越界的 \code{i}，因此返回类型不需要 \code{Option A}。代价是程序变换列表长度后，相关索引也可能需要 transport。

\section{无效状态不可表示，还是允许后证明}

建模有两种基本路线：
\begin{description}
  \item[精确类型] 把长度、边界或合法性编码进类型，使无效状态不可构造；
  \item[普通数据加不变量] 使用简单结构表示状态，再单独定义 \code{Inv s : Prop} 并证明转换保持它。
\end{description}
前者让部分错误在类型检查阶段消失，后者通常更容易编程、重写和对接现有实现。选择标准不是“依赖类型越多越高级”，而是证明成本与接口收益。

\section{structure：命名字段的乘积类型}

\begin{lstlisting}
structure StreamState where
  seen : Nat
  kept : List Nat
deriving Repr

def StreamState.push (s : StreamState) (x : Nat) : StreamState :=
  { seen := s.seen + 1, kept := x :: s.kept }

def StreamState.Inv (s : StreamState) : Prop :=
  s.kept.length <= s.seen
\end{lstlisting}
结构投影 \code{s.seen} 和 \code{s.kept} 是普通函数。更新语法产生新值，不会原地修改旧状态；这使状态转换适合直接写成数学函数。

\section{一个保持不变量的完整小定理}

\begin{lstlisting}
theorem StreamState.push_preserves
    (s : StreamState) (x : Nat)
    (h : StreamState.Inv s) :
    StreamState.Inv (s.push x) := by
  simpa [StreamState.Inv, StreamState.push] using
    Nat.succ_le_succ h
\end{lstlisting}
展开后，目标是
\[
\operatorname{length}(x::s.\mathrm{kept})
\le s.\mathrm{seen}+1,
\]
即从 $\operatorname{length}(s.\mathrm{kept})\le s.\mathrm{seen}$ 两边同时加一。这个例子展示了算法形式化的基本结构：状态、转换、不变量、保持性。

\section{Subtype：值与性质打包}

子类型
\[
\{x:A\mid P(x)\}
\]
包含一个值和它满足 $P$ 的证明。例如：
\begin{lstlisting}
def PositiveNat := {n : Nat // 0 < n}

def onePositive : PositiveNat :=
  Subtype.mk 1 (by decide)

#check onePositive.val
#check onePositive.property
\end{lstlisting}
使用 subtype 时，值层计算常可忽略证明字段，但等式和转换有时需要处理 coercion。若大量时间都花在搬运证明而不是目标性质上，应重新评估接口。

\section{Sigma 类型与存在命题}

依赖对 \code{Sigma B} 保存一个索引 $a:A$ 和一个类型为 $B(a)$ 的数据。它位于 \code{Type}，通常用于程序数据；\code{Exists} 位于 \code{Prop}，主要表达逻辑存在，证明在计算提取中的行为不同。

可以粗略记成：
\[
\Sigma_{a:A} B(a) \quad\text{保存可计算数据},\qquad
\exists a:A,\ P(a) \quad\text{声明命题成立}.
\]
实际使用时还要检查 universe、proof irrelevance 与是否需要从证明中提取 witness。

\section{等式与 dependent transport}

若 $h:a=b$，普通重写把关于 $a$ 的表达式替换为关于 $b$ 的表达式。依赖类型中，$B(a)$ 与 $B(b)$ 可能是表面不同的类型，需要沿 $h$ 搬运数据。这是 \code{Eq.rec}、\code{cast}、\code{subst} 等机制背后的问题。

\begin{lstlisting}
example (a b : Nat) (h : a = b) (P : Nat -> Prop)
    (ha : P a) : P b := by
  cases h
  exact ha
\end{lstlisting}
对 $h$ 做 cases 后只剩反身情形 $a=a$，目标退化为已有的 \code{ha}。

\begin{warningbox}
数学上“长度相等”不代表 Lean 会自动把 \code{Fin xs.length} 当成 \code{Fin ys.length}。依赖索引属于类型的一部分。先寻找能避免 transport 的高层 lemma；只有确实需要时再显式搬运。
\end{warningbox}

\section{Typeclass：由类型驱动的接口搜索}

Typeclass 把操作和定律接口化，实例参数用方括号表示：
\begin{lstlisting}
class HasCost (A : Type) where
  cost : A -> Nat

def totalCost [HasCost A] (xs : List A) : Nat :=
  (xs.map HasCost.cost).sum
\end{lstlisting}
当调用 \code{totalCost} 时，elaborator 需要找到 \code{HasCost A} 实例。代数结构、序、可判定命题、有限类型等 mathlib 接口大量依赖这种搜索。

实例找不到时，应检查：目标类型是否真的是预期类型、实例所在模块是否已导入、是否存在局部歧义、参数有没有被 coercion 改变。不要把所有问题都归因于“mathlib 没这个定理”。

\section{可判定性与计算}

要在程序中对命题做 \code{if} 分支，Lean 需要 \code{Decidable p}，即能计算判断 $p$ 真假的接口。有限等式、自然数序关系等通常已有实例。\code{by decide} 可以关闭可计算的可判定命题，但它证明的是当前形式表达式，不会自动选择正确建模。

\begin{aibox}
自动形式化系统经常生成“类型上很强、语义上很怪”的 statement，例如把非法输入排除得过多，使 theorem 变得平凡。依赖类型可以增强规范，也可以隐藏任务偷换；必须把自然语言输入域、Lean 类型和排除条件逐项对照。
\end{aibox}

\begin{checkpointbox}
\begin{enumerate}
  \item 用 \code{Fin n} 解释安全索引为何不会越界；
  \item 比较“精确类型”和“结构加不变量”两种状态建模方式；
  \item 编译 \code{StreamState.push\_preserves} 并逐步展开目标；
  \item 构造一个 subtype 值，分别访问值与证明字段；
  \item 说明 dependent transport 和普通等式重写有何不同；
  \item 诊断一个 typeclass instance synthesis 失败的最小例子。
\end{enumerate}
\end{checkpointbox}

\chapter{mathlib 检索、自动化与证明调试}

\section{先读目标，再搜名字}

mathlib 很大，靠猜定理名只能偶尔成功。稳定流程是：
\begin{enumerate}
  \item 读当前 local context 与 target；
  \item 确认对象的具体类型、命名空间和 coercion；
  \item 用 \code{\#check}、文档和源代码找候选定理；
  \item 在最小例子中测试候选的参数方向；
  \item 最后才把它接回原证明。
\end{enumerate}

\begin{lstlisting}
#check Nat.add_comm
#check Nat.succ_le_succ
#check List.reverse_reverse
#print Nat.add_comm
\end{lstlisting}
\code{\#check} 看接口，\code{\#print} 看声明展开和依赖。阅读库代码时优先寻找同一命名空间中相邻定理，因为命名通常体现对象与操作。

\section{目标形状比自然语言关键词更可靠}

自然语言说“单调”，Lean 目标可能是函数单调、严格单调、集合包含或序同态；说“有限”，可能对应 \code{Finset}、\code{Fintype}、\code{Set.Finite} 或有限维空间。先确定形式对象，搜索才不会跨错抽象层。

例如看到
\[
a+1\le b+1
\]
应先识别它是自然数序上的 successor 单调，而不是搜索任何含“加法不等式”的定理。

\section{建议型搜索 tactic}

mathlib 提供能建议候选脚本的工具：
\begin{center}
\begin{tabularx}{0.94\textwidth}{>{\ttfamily}lX}
\toprule
\normalfont 工具 & 用途 \\
\midrule
exact? & 搜索可直接匹配当前目标的项 \\
apply? & 搜索应用后能产生可处理子目标的定理 \\
rw? & 建议一条可能推进目标的重写 \\
simp? & 展示建议的简化规则集合 \\
aesop? & 对构造子、假设和安全规则进行结构化搜索 \\
\#check & 确认候选定理的完整类型 \\
\#print & 查看定义或定理的展开信息 \\
\bottomrule
\end{tabularx}
\end{center}

问号版本的价值在于返回一段可读建议。把它改写成稳定、较小的显式脚本后，应重新读懂每条依赖；不要长期把宽泛搜索当作唯一证明说明。

\section{常见领域自动化}

不同 tactic 解决不同理论：
\begin{description}
  \item[\code{norm\_num}] 具体数字、半环/环中的数值归约；
  \item[\code{ring} / \code{ring\_nf}] 交换半环或环中的多项式恒等式；
  \item[\code{linarith}] 由线性等式与不等式组合推出目标；
  \item[\code{nlinarith}] 包含有限非线性多项式关系的算术推理；
  \item[\code{omega}] Presburger 算术，尤其自然数和整数的线性关系；
  \item[\code{decide}] 对有可判定实例且可计算的封闭命题求证；
  \item[\code{native\_decide}] 借助更高效计算检查可判定命题，使用边界需按项目规范记录。
\end{description}

\begin{lstlisting}
example : (17 : Nat) < 23 := by
  norm_num

example (a b : Int) (h1 : a <= b) : a + 3 <= b + 3 := by
  omega

example (x y : Int) : (x + y)^2 = x^2 + 2*x*y + y^2 := by
  ring
\end{lstlisting}
自动化失败时先检查问题是否属于它支持的理论。例如 \code{linarith} 不负责列表归纳，\code{ring} 不知道任意非环函数的语义。

\section{控制 simp 的作用域}

\code{simp} 依赖一个不断扩展的规则集。学习与稳定证明中常用：
\begin{lstlisting}
simp [myDefinition]
simp only [List.length_cons, Nat.succ_le_succ_iff]
simpa [myDefinition] using h
\end{lstlisting}
\code{simp only} 让依赖更明确，但可能冗长；普通 \code{simp} 更易维护，但受全局 simp 集变化影响。关键证明可用 \code{simp?} 查看实际规则，再决定保留哪种粒度。

不要随意把双向都可能使用的定理标记为 simp，否则可能造成循环、性能下降或不可预测的规范形。理想 simp lemma 把表达式推向明确的正常形。

\section{类型不匹配的诊断}

“看起来一样”却无法应用，常见原因包括：
\begin{itemize}
  \item \code{Nat} 与 \code{Int} 之间存在 coercion；
  \item 一个对象是 \code{List A}，另一个是 \code{Finset A}；
  \item 隐式 universe 或类型参数不同；
  \item 目标经过 reducible definition 展开后才匹配；
  \item 定理需要 typeclass 实例，但当前没有或不唯一；
  \item 等式方向、函数参数顺序或命名空间不对。
\end{itemize}

可临时使用更详细打印：
\begin{lstlisting}
set_option pp.all true in
#check Nat.add_comm
\end{lstlisting}
详细输出很嘈杂，只用于定位隐式结构；问题解决后应恢复可读形式。

\section{最小失败例子}

可靠调试流程如下：
\begin{center}
\begin{tikzpicture}[
  node distance=5mm,
  box/.style={draw,rounded corners=1mm,align=center,minimum width=5.6cm,minimum height=7mm},
  arr/.style={-{Stealth[length=2mm]},thick}
]
\node[box] (a) {复制失败 theorem 到独立文件};
\node[box,below=of a] (b) {删除无关定义与 imports};
\node[box,below=of b] (c) {固定一个 proof state};
\node[box,below=of c] (d) {检查 statement、类型和候选 lemma};
\node[box,below=of d] (e) {手写一小步，再引入自动化};
\node[box,below=of e] (f) {回填原文件并运行 lake build};
\draw[arr] (a)--(b); \draw[arr] (b)--(c); \draw[arr] (c)--(d);
\draw[arr] (d)--(e); \draw[arr] (e)--(f);
\end{tikzpicture}
\end{center}

记录至少包括：完整错误、当前目标、上下文、尝试过的定理及其 \code{\#check} 输出、最终原因分类。这样的失败轨迹以后可以成为检索数据或 AI proof repair 的训练/评测样本。

\section{不要用过宽 import 掩盖依赖}

\code{import Mathlib} 很适合入门和探索；稳定模块可逐步缩小 imports，以明确依赖和改善构建性能。但过早手工优化 imports 会打断学习，合理顺序是：先让 theorem 在固定版本编译，再用工具或构建错误识别真实依赖。

\begin{warningbox}
禁止用 \code{sorry}、额外 \code{axiom} 或把目标改弱来伪装修复。学习时可以暂时插入 \code{sorry} 隔离后续错误，但提交前必须在状态记录中显式统计，并区分“文件可编译”和“证明无占位”。
\end{warningbox}

\begin{aibox}
一个有价值的 proof-repair benchmark 不应只记录最终成功脚本，还应保存 Lean/toolchain 版本、原始错误、proof state、允许的依赖和语义不变约束。否则模型可能通过修改 statement 或引入强公理获得虚假成功。
\end{aibox}

\begin{checkpointbox}
\begin{enumerate}
  \item 从目标类型出发，在 mathlib 文档中找到一个所需定理；
  \item 分别用 \code{norm\_num}、\code{omega} 和 \code{ring} 解决适合它们的问题；
  \item 用 \code{simp?} 查看建议，并改写成一段你能解释的脚本；
  \item 把一个失败 theorem 缩成 20 行以内的最小例子；
  \item 说明 \code{sorry}、\code{axiom}、测试通过和 kernel-checked theorem 的证据差别。
\end{enumerate}
\end{checkpointbox}

\chapter{算法不变量、有限随机性与证明分层}

\section{先把随机性变成显式输入}

普通程序常在函数内部调用随机数生成器，这不利于数学推理。形式化时可以把所有随机选择看成显式输入 $r$：
\[
\operatorname{run}:\mathrm{Input}\to\mathrm{RandomTape}\to\mathrm{Output}.
\]
固定 $r$ 后，\code{run x r} 是确定函数；概率来自对 random tape 的分布。这样可以把问题拆成两层：
\begin{enumerate}
  \item 对每条随机带证明状态安全、输出类型正确等确定性性质；
  \item 对随机带集合计数或使用概率测度，证明成功事件占多大比例。
\end{enumerate}

\begin{intuitionbox}
随机算法不是“不可预测的代码”，而是“确定程序加一个按指定分布采样的额外输入”。这个视角让执行语义与概率分析分离，也便于用有限枚举、Lean 和外部测试交叉验证。
\end{intuitionbox}

\section{状态机、不变量与后置条件}

流算法可写成：
\[
s_0=\operatorname{init},\qquad
s_{t+1}=\operatorname{step}(s_t,x_{t+1},r_{t+1}),\qquad
y=\operatorname{query}(s_n).
\]
至少区分三类性质：
\begin{description}
  \item[表示不变量] 状态字段长度、索引、非负性和内部一致性；
  \item[语义不变量] 状态和已处理输入前缀之间的关系；
  \item[最终保证] 在终止或查询时由不变量推出输出正确性/误差界。
\end{description}

上一章的 \code{StreamState.Inv} 只保证列表长度不超过已见元素数，它是表示不变量，不说明保留的元素具有正确抽样分布。

\begin{warningbox}
“状态大小不超过 $k$”不是 reservoir sampling 正确性的全部。真正分布性结论是处理 $n$ 个元素后，每个元素以 $k/n$ 概率进入样本（在相应条件下），还需要形式化随机选择及其分布。
\end{warningbox}

\section{有限均匀样本空间}

入门时可以绕开测度论，令有限类型 $\Omega$ 表示所有随机带，均匀概率定义为
\[
\Pr[E]=\frac{|\{\omega\in\Omega:E(\omega)\}|}{|\Omega|},
\qquad |\Omega|>0.
\]
Lean 中可先定义事件计数：
\begin{lstlisting}
def eventCount [Fintype Omega]
    (E : Omega -> Prop) [DecidablePred E] : Nat :=
  (Finset.univ.filter E).card

def uniformProb [Fintype Omega]
    (E : Omega -> Prop) [DecidablePred E] : Rat :=
  (eventCount E : Rat) / (Fintype.card Omega : Rat)
\end{lstlisting}
若 $\Omega$ 为空，分母为零，必须在定理假设中排除或改用带非空证明的样本空间结构。定义“能通过类型检查”不等于数学接口已完整。

\section{indicator 与期望}

事件 $E$ 的指示变量
\[
\mathbf{1}_E(\omega)=
\begin{cases}
1,&E(\omega),\\
0,&\text{otherwise}
\end{cases}
\]
满足 $\mathbb{E}[\mathbf{1}_E]=\Pr[E]$。有限均匀空间中，这一事实归结为：对所有 $\omega$ 求和时，每个满足 $E$ 的元素贡献 1，其余贡献 0，因此总和恰为事件基数。

形式化路线可以拆成：
\begin{enumerate}
  \item 定义可判定事件和 indicator；
  \item 证明 indicator 的有限和等于 filter 后的 card；
  \item 除以非零的 $|\Omega|$；
  \item 再把结果包装成期望/概率接口。
\end{enumerate}
这比直接挑战一般概率库中的积分定理更适合作为第一题。

\section{有限 union bound 的计数骨架}

对事件 $E_1,\ldots,E_m$，union bound 为
\[
\Pr\!\left[\bigcup_{i=1}^m E_i\right]
\le \sum_{i=1}^m \Pr[E_i].
\]
有限均匀空间下只需证明集合基数不等式
\[
\left|\bigcup_i E_i\right|\le\sum_i |E_i|,
\]
再除以 $|\Omega|$。重叠元素在右边可能被重复计数，所以得到上界。Lean 实现时要选择 \code{Finset}、\code{Set} 或有限索引族，并处理 membership、decidability 与 coercion。

\section{随机算法保证的六层拆分}

\begin{center}
\begin{tabularx}{0.96\textwidth}{>{\bfseries}lX X}
\toprule
层 & 典型命题 & 合适证据 \\
\midrule
接口 & 输入可解析、输出类型正确 & 类型检查、编译 \\
状态安全 & 索引合法、不变量保持 & Lean 归纳定理 \\
逐随机带正确性 & 对所有 $r$ 都满足确定性关系 & Lean theorem \\
分布正确性 & random tape 确实服从所需分布 & 有限计数或概率库定理 \\
概率保证 & 失败概率不超过 $\delta$ & 期望、尾界、union bound 的形式证明 \\
资源保证 & 时间、空间、随机位成本 & 成本语义、复杂度证明或审计分析 \\
\bottomrule
\end{tabularx}
\end{center}

一个项目可以只形式化其中一部分，但必须标明其余层是测试、纸笔证明、外部假设还是完全未处理。

\section{Count-Min 的分层形式化示例}

Count-Min Sketch 的端到端结论可以先拆成：
\begin{enumerate}
  \item 数组更新保持所有计数器非负；
  \item 目标项每行对应桶包含 $f_i$，所以逐行估计不小于 $f_i$；
  \item 碰撞噪声是其他频率的加权和；
  \item 通用哈希给出碰撞概率，推出噪声期望界；
  \item Markov 与独立重复给出高概率误差界；
  \item 最小值查询和参数 $w,d$ 合成最终 theorem。
\end{enumerate}
第一、二步基本不含概率，适合先做；后三步需要哈希族和有限概率接口。完整形式化不是“一条巨大 theorem”，而是一条依赖图。

\section{实现与规范的连接}

Lean 中容易证明的纯函数可能与 Python/C++ 高性能实现不同。至少有三种连接策略：
\begin{enumerate}
  \item 以 Lean 定义为可执行 reference implementation，外部实现做 differential testing；
  \item 为外部语言建立语义模型，再证明实现 refinement；
  \item 只证明抽象算法 theorem，同时明确工程实现尚未形式化。
\end{enumerate}
第一种最适合个人学习项目。它不能自动证明外部实现正确，但能给测试提供可信 oracle。

\begin{aibox}
AI 生成的随机算法候选应先通过分层 verifier：接口检查、小规模枚举、hidden random tests、确定性不变量证明、概率保证审查和复杂度审查。Lean 最适合其中可精确定义的层，不应被当作覆盖所有研究判断的万能标签。
\end{aibox}

\section{建议的第一个随机算法形式化题}

选择如下梯度：
\begin{enumerate}
  \item 有限集合上 indicator 的和等于事件基数；
  \item 均匀概率总和为 1（显式假设样本空间非空）；
  \item 两个或有限多个事件的 union bound；
  \item 一个抽样状态的大小不变量；
  \item 简化抽样器的无偏性；
  \item 最后才进入完整 reservoir sampling、Chernoff 或 Count-Min。
\end{enumerate}

\begin{checkpointbox}
\begin{enumerate}
  \item 把随机算法写成输入和 random tape 上的确定函数；
  \item 区分表示不变量、语义不变量和分布正确性；
  \item 在有限均匀空间中推导 indicator 期望等于事件概率；
  \item 写出 union bound 的集合计数证明骨架；
  \item 把 Count-Min 端到端保证拆成至少五条 lemma；
  \item 为一个只完成部分形式化的项目写清证据边界。
\end{enumerate}
\end{checkpointbox}

\chapter{Lean 在 AI4Math 与 AI4TCS 中的证明工程}

\section{Lean 在闭环中的位置}

一个典型闭环为：
\begin{center}
\begin{tikzpicture}[
  node distance=6mm,
  box/.style={draw,rounded corners=1mm,align=center,minimum width=4.8cm,minimum height=8mm},
  arr/.style={-{Stealth[length=2mm]},thick}
]
\node[box] (nl) {自然语言问题与背景};
\node[box,below=of nl] (stmt) {形式 statement 与定义};
\node[box,below=of stmt] (gen) {模型/搜索器生成 proof 或 program};
\node[box,below=of gen] (lean) {Lean elaboration 与 kernel 检查};
\node[box,below=of lean] (fb) {proof state、错误与反例反馈};
\node[box,below=of fb] (human) {人类审查语义、假设、意义与复杂度};
\draw[arr] (nl)--(stmt); \draw[arr] (stmt)--(gen); \draw[arr] (gen)--(lean);
\draw[arr] (lean)--(fb); \draw[arr] (fb)--(human);
\draw[arr] (fb.east) .. controls +(2.1,0) and +(2.1,0) .. (gen.east);
\end{tikzpicture}
\end{center}

Lean 是强 verifier，不是自动保证研究全流程正确的神谕。它精确回答“这个 proof term 是否证明了这个形式 statement”，其余问题需要额外证据。

\section{AI4Math 的四类任务}

\begin{description}
  \item[定理证明] statement 已给定，在固定环境中搜索 kernel 可接受证明；
  \item[自动形式化] 把自然语言、公式和上下文翻译成 Lean 定义与 statement；
  \item[proof repair] 给定失败代码和环境，修复证明而保持 statement 语义不变；
  \item[猜想与对象发现] 生成候选定理、反例、构造或数学对象，再经形式/数值工具筛选。
\end{description}
这四类任务的输入、成功条件和风险不同。定理证明 benchmark 上的成功率不能直接代表自动形式化忠实度，更不能代表发现了有价值的新数学。

\section{statement fidelity 是独立任务}

自然语言到 Lean 至少要核对：
\begin{enumerate}
  \item 量词范围：所有对象、存在对象还是某个固定对象；
  \item 输入域：自然数、整数、实数、有限图还是任意图；
  \item 前提：非空、独立、连通、可测、有限等条件是否遗漏；
  \item 等价概念：自然语言术语是否与库定义同义；
  \item 结论强度：严格不等式、近似界、概率至少多少；
  \item 退化情形：是否通过加入过强假设或空类型让命题平凡成立。
\end{enumerate}

推荐维护三列表：自然语言片段、数学形式、Lean 片段，并对每一处量词和假设做对齐。kernel 无法替你完成这张表。

\section{proof search 的状态、动作与反馈}

形式证明搜索可以建模为：
\begin{description}
  \item[state] 当前 goals、local context、已导入环境与选项；
  \item[action] tactic、term、lemma application 或 statement decomposition；
  \item[transition] Lean 执行动作后产生新 goals 或错误；
  \item[reward/feedback] 是否编译、目标数量变化、错误类型、证明长度与资源成本；
  \item[terminal] 无 goals 且无占位，kernel 接受最终声明。
\end{description}

仅以“编译通过”为奖励会鼓励系统寻找捷径，例如修改 statement、加入不允许的公理或利用数据泄漏。评测必须冻结目标和允许环境，并审计依赖。

\section{proof repair 的不变约束}

一个真正的 proof repair 任务至少固定：
\begin{itemize}
  \item theorem statement；
  \item Lean/mathlib 版本；
  \item 允许修改的代码范围；
  \item 禁止新增的公理、\code{sorry} 和不透明外部假设；
  \item 资源预算与是否允许网络检索；
  \item 最终要运行的干净构建命令。
\end{itemize}
若允许修改 statement，任务就变成联合的 statement repair，需要单独评估语义是否保留。

\section{AI4TCS 中的算法验证}

候选算法通常需要多种 verifier：
\begin{enumerate}
  \item parser/type checker 确认接口；
  \item 小规模穷举寻找功能反例；
  \item 随机与 OOD 测试检查经验泛化；
  \item SMT/SAT 检查有限域或约束性质；
  \item Lean 证明不变量、refinement 或数学 theorem；
  \item 人类分析时间、空间、随机位和渐近复杂度。
\end{enumerate}
Lean 适合组合已精确定义的性质，但性能结果通常需要成本模型；“函数定义可执行”并不自动给出 $O(n\log n)$。

\section{证据登记表}

\begin{center}
\begin{tabularx}{0.96\textwidth}{>{\bfseries}lX X}
\toprule
证据 & 可以支持 & 不能单独支持 \\
\midrule
Lean kernel 接受 & 当前声明在当前环境中可证明 & 翻译忠实、原创、有用、实现高效 \\
小规模穷举 & 范围内无反例或找到反例 & 任意规模 correctness \\
随机测试 & 指定分布与预算下表现 & worst-case guarantee \\
SMT/SAT unsat & 编码约束下没有反例 & 编码忠实、无限域结论 \\
基准分数 & 固定 benchmark 上的经验能力 & 开放世界研究能力 \\
复杂度证明 & 成本模型下的资源上界 & 常数、工程性能、模型外实现 \\
\bottomrule
\end{tabularx}
\end{center}

\section{数据集与污染问题}

形式定理有精确字符串和依赖图，近重复与泄漏尤其容易发生。拆分 benchmark 时应考虑：
\begin{itemize}
  \item 同一定理的重命名、等价改写和自动生成变体；
  \item theorem 之间的依赖边，测试定理是否已作为训练 lemma 出现；
  \item 仓库时间切分与版本历史；
  \item 自然语言题与形式 statement 的跨语料对应；
  \item 搜索时访问的文档、检索库和缓存。
\end{itemize}
随机按 theorem 行拆分通常不足以排除结构泄漏。

\section{失败轨迹也是研究产物}

成功 proof 只展示一条路径；失败轨迹包含模型真正缺失的信息。建议保存：
\begin{enumerate}
  \item 原始 statement 与环境 hash；
  \item 每一步 proof state 和候选动作；
  \item parser、elaboration、timeout、unknown lemma 等错误分类；
  \item 最小失败例子与人工修复理由；
  \item 是否修改定义、假设或导入；
  \item 最终 proof 的依赖和构建日志摘要。
\end{enumerate}
这些信息可用于错误分析、proof repair、检索改进与可复现实验，而不只是“成功/失败”一位标签。

\section{个人可完成的最小闭环}

选择一个 20--50 行的有限对象 theorem：
\begin{enumerate}
  \item 写自然语言、数学式和 Lean statement 对照；
  \item 手写 baseline proof；
  \item 人为制造 5 类局部错误；
  \item 让模型或搜索工具尝试修复，并保存每次错误反馈；
  \item 在隐藏的同构变体上测试；
  \item 运行干净 \code{lake build}；
  \item 报告成功率、失败类型、证明依赖与语义审查结果。
\end{enumerate}

\begin{warningbox}
不要以“Lean 编译成功”概括整个项目。准确表述应指出：哪个 statement、哪个 toolchain、是否无 \code{sorry}/额外公理、自然语言对齐如何检查、实现和复杂度是否在形式化范围内。
\end{warningbox}

\begin{checkpointbox}
\begin{enumerate}
  \item 区分 theorem proving、autoformalization、proof repair 与 conjecture discovery；
  \item 为一个自然语言题制作 statement fidelity 对照表；
  \item 把 proof search 写成 state/action/transition/terminal；
  \item 为 proof repair 任务写出不可修改约束；
  \item 设计一个 Lean、穷举、随机测试和人工审查分工的 verifier；
  \item 解释为什么随机 theorem split 可能产生依赖泄漏。
\end{enumerate}
\end{checkpointbox}

\chapter{综合习题与参考答案}

\section{习题}

\begin{enumerate}[label=\textbf{L\arabic*.}]
  \item 画出 Lean 源码从 parser/elaborator 到 kernel 的处理路径，并说明 tactic 为什么不必属于可信内核。
  \item 建立最小 Lake 项目，写出需要提交以复现构建的文件，以及不应提交的缓存目录。
  \item 定义 \code{triple : Nat -> Nat}，用 \code{\#eval} 计算一个例子，再证明 \code{triple n = n+n+n}。
  \item 分别用 proof term 与 tactic 证明 \code{p -> q -> And p q}。
  \item 证明 \code{Or p q -> Or q p}，解释为什么必须做两个分支。
  \item 对列表证明 \code{List.reverse (List.reverse xs) = xs}。若直接搜索 mathlib 引理，记录其完整类型；若手工证明，写出归纳结构。
  \item 定义一棵二叉树和节点计数函数，证明镜像保持节点数。
  \item 给出一个归纳假设太弱的例子，写出加强后的 theorem statement。
  \item 用 \code{Fin xs.length} 定义安全索引，说明为什么空列表上无法构造索引参数。
  \item 对 \code{StreamState.push\_preserves} 逐步写出展开前目标、展开后目标和所用自然数引理。
  \item 构造一个 mathlib theorem 搜索失败的最小例子，按“定义/导入/类型/实例/引理”分类根因。
  \item 在有限均匀样本空间中证明 indicator 的期望等于事件概率，列出分母非零假设。
  \item 把 reservoir sampling 的正确性至少拆成状态安全、随机选择分布和最终边缘概率三层。
  \item 为一个自然语言定理写自然语言、数学式、Lean statement 三列对照，找出至少两个可能的语义丢失点。
  \item 设计一个 proof-repair 小实验：固定 statement 和 toolchain，定义允许修改范围、五类错误、指标与 hidden 变体。
\end{enumerate}

\section{参考答案与提示}

\paragraph{L1.}
路径为源码/tactic $\to$ parser 与 elaborator $\to$ 带完整隐式信息的核心表达式/证明项 $\to$ kernel 类型检查。Tactic 只负责构造候选证明项；只要最终项由 kernel 复查，tactic 自身出错通常导致拒绝或产生错误候选，而不是直接让错误 theorem 被接受。

\paragraph{L2.}
至少提交 \code{lean-toolchain}、Lake 配置、manifest、所有项目源码和构建说明。\code{.lake/}、编辑器状态和本机缓存不作为源代码提交。复现命令应从干净 checkout 运行 \code{lake build}。

\paragraph{L3.}
\begin{lstlisting}
def triple (n : Nat) : Nat := n + n + n

#eval triple 4

example (n : Nat) : triple n = n + n + n := by
  rfl
\end{lstlisting}
\code{rfl} 成功是因为展开定义后两边相同；\code{\#eval} 只展示输入 4 的结果。

\paragraph{L4.}
直接项为 \code{fun hp hq => And.intro hp hq}。Tactic 版先 \code{intro hp hq}，再 \code{exact And.intro hp hq}，或者 \code{constructor <;> assumption}。应能解释最终项仍是一连串函数与构造子。

\paragraph{L5.}
对输入证明 \code{h : Or p q} 做 \code{cases h}。左分支携带 \code{hp : p}，输出 \code{Or.inr hp}；右分支携带 \code{hq : q}，输出 \code{Or.inl hq}。析取值可能由任一构造子产生，漏掉分支就不是总证明。

\paragraph{L6.}
mathlib 已有 \code{List.reverse\_reverse}，应先 \code{\#check List.reverse\_reverse} 并按当前版本的参数应用。手工证明对 \code{xs} 归纳，\code{cons} 分支需要 reverse 与 append 的关系；这也说明复用库 lemma 比从头展开更合适。

\paragraph{L7.}
按 \code{empty/node} 归纳。空树直接 \code{rfl}；节点分支使用左右子树两个归纳假设，再用自然数加法交换律/结合律处理镜像交换后的顺序。正文 \code{BinTree.size\_mirror} 给出参考骨架。

\paragraph{L8.}
尾递归求和若只证明初始累加器为 0，归纳分支通常需要对新的累加器 $acc+x$ 使用假设。应加强为“对任意 $acc$，尾递归结果等于 $acc+\mathrm{sum}(xs)$”，并在归纳时保持 $acc$ 一般。

\paragraph{L9.}
\code{safeGet (xs) (i : Fin xs.length)} 的 \code{i} 包含一个自然数值与小于 \code{xs.length} 的证明。当 \code{xs=[]} 时需要构造 \code{Fin 0}，即给出小于 0 的自然数，这是不可能的。

\paragraph{L10.}
展开前目标是 \code{Inv (s.push x)}；展开后是 \code{(x::s.kept).length <= s.seen+1}，化简为 \code{s.kept.length+1 <= s.seen+1}。由假设和 \code{Nat.succ\_le\_succ} 得到，\code{simpa} 只负责连接表达形式。

\paragraph{L11.}
合格记录要包含最小源码、toolchain、完整错误、proof state、候选定理的 \code{\#check} 和最终分类。若换一个 import 就成功，根因是依赖；若 \code{@theorem} 显示参数类型不同，根因是类型/隐式参数；若缺 \code{DecidableEq A}，根因是实例或缺少假设。

\paragraph{L12.}
对 indicator 在 $\Omega$ 上求和，恰得到满足 $E$ 的元素个数。两边除以 $|\Omega|$ 得结论。必须假设 $\Omega$ 有限且非空，并明确均匀分布；事件需要可判定以便有限计算。

\paragraph{L13.}
状态安全包括样本大小不超过 $k$、索引合法；随机选择分布说明第 $n$ 步替换位置的抽样方式；最终边缘概率通过对 $n$ 归纳证明每个已见元素被保留的概率为 $\min(1,k/n)$。只完成第一层不能声称 reservoir sampling 分布正确。

\paragraph{L14.}
重点检查量词、输入域、非空/有限/独立等假设、严格或非严格关系，以及退化输入。例如把“任意连通图”误译成“任意图”，或额外假设候选解已经最优，都会改变语义，即使 Lean 能证明新 statement。

\paragraph{L15.}
固定 theorem 文本、imports 和 toolchain；只允许改 proof body。错误可包括错 lemma 名、等式方向、缺失分支、归纳变量错误和实例缺失。报告修复成功率、平均尝试数、错误分类、proof 长度、是否引入禁止依赖，并在变量重命名/同构改写的 hidden theorem 上复测。


\begin{thebibliography}{9}
\bibitem{tpil}
Jeremy Avigad, Leonardo de Moura, Soonho Kong, and contributors.
\newblock \emph{Theorem Proving in Lean 4}.
\newblock \url{https://lean-lang.org/theorem_proving_in_lean4/}.

\bibitem{fpl}
David Thrane Christiansen.
\newblock \emph{Functional Programming in Lean}.
\newblock \url{https://lean-lang.org/functional_programming_in_lean/}.

\bibitem{mil}
Jeremy Avigad, Patrick Massot, and Floris van Doorn.
\newblock \emph{Mathematics in Lean}.
\newblock \url{https://leanprover-community.github.io/mathematics_in_lean/}.

\bibitem{reference}
The Lean Team.
\newblock \emph{Lean Language Reference}.
\newblock \url{https://lean-lang.org/doc/reference/latest/}.

\bibitem{mathlib}
The mathlib Community.
\newblock \emph{mathlib4 Documentation}.
\newblock \url{https://leanprover-community.github.io/mathlib4_docs/}.
\end{thebibliography}

\end{document}
